\section{AMC 10 10A 2003}
        \begin{enumerate}
        	\item (\textit{AMC 10 10A 2003}) What is the difference between the sum of the first $2003$ even counting numbers and the sum of the first $2003$ odd counting numbers? 


 $\mathrm{(A) \ } 0\qquad \mathrm{(B) \ } 1\qquad \mathrm{(C) \ } 2\qquad \mathrm{(D) \ } 2003\qquad \mathrm{(E) \ } 4006$

 \textbf{Answer: }D



\textbf{Solution 1}

The first $2003$ even counting numbers are $2,4,6,...,4006$. 

 The first $2003$ odd counting numbers are $1,3,5,...,4005$. 

 Thus, the problem is asking for the value of $(2+4+6+...+4006)-(1+3+5+...+4005)$. 

 $(2+4+6+...+4006)-(1+3+5+...+4005) = (2-1)+(4-3)+(6-5)+...+(4006-4005)$
$= 1+1+1+...+1 = \boxed{\mathrm{(D)}\ 2003}$

 The answer is 2003

 



\textbf{Solution 2}

Using the sum of an \href{https://artofproblemsolving.com/wiki/index.php?title=Arithmetic\_progression}{arithmetic progression} formula, we can write this as $\frac{2003}{2}(2 + 4006) - \frac{2003}{2}(1 + 4005) = \frac{2003}{2} \cdot 2 = \boxed{\mathrm{(D)}\ 2003}$.

 


 



\textbf{Solution 3}

The formula for the sum of the first $n$ even numbers, is $S_E=n^{2}+n$, (E standing for even).

 Sum of first $n$ odd numbers, is $S_O=n^{2}$, (O standing for odd).

 Knowing this, plug $2003$ for $n$, 

 $S_E-S_O= (2003^{2}+2003)-(2003^{2})=2003 \Rightarrow$ $\boxed{\mathrm{(D)}\ 2003}$.

 



\textbf{Solution 4}

In the case that we don't know if $0$ is considered an even number, we note that it doesn't matter! The sum of odd numbers is $O=1+3+5+...+4005$. And the sum of even numbers is either $E_1=0+2+4...+4004$ or $E_2=2+4+6+...+4006$. When compared to the sum of odd numbers, we see that each of the $n$th term in the series of even numbers differ by $1$. For example, take series $O$ and $E_1$. The first terms are $1$ and $0$. Their difference is $|1-0|=1$. Similarly, take take series $O$ and $E_2$. The first terms are $1$ and $2$. Their difference is $|1-2|=1$. Since there are $2003$ terms in each set, the answer $\boxed{\mathrm{(D)}\ 2003}$.

 



\textbf{Solution 5 (Fastest method)}

We can pair each term of the sums - the first even number with the first odd number, the second with the second, and so forth. Then, there are 2003 pairs with a difference of 1 in each pair - 2-1 is 1, 4-3 is 1, 6-5 is 1, and so on. Then, the solution is $1 \cdot 2003$, and the answer is $\boxed{\text{(D) }2003}$.

 <3

 


	\item (\textit{AMC 10 10A 2003}) Members of the Rockham Soccer League buy socks and T-shirts. Socks cost $4 per pair and each T-shirt costs $5 more than a pair of socks. Each member needs one pair of socks and a shirt for home games and another pair of socks and a shirt for away games. If the total cost is $2366, how many members are in the League? 


 $\mathrm{(A) \ } 77\qquad \mathrm{(B) \ } 91\qquad \mathrm{(C) \ } 143\qquad \mathrm{(D) \ } 182\qquad \mathrm{(E) \ } 286$

 \textbf{Answer: }B



\textbf{Solution}

Since T-shirts cost $5$ dollars more than a pair of socks, T-shirts cost $5+4=9$ dollars. 

 Since each member needs $2$ pairs of socks and $2$ T-shirts, the total cost for $1$ member is $2(4+9)=26$ dollars. 

 Since $2366$ dollars was the cost for the club, and $26$ was the cost per member, the number of members in the League is $2366\div 26=\boxed{\mathrm{(B)}\ 91}$. 

 


	\item (\textit{AMC 10 10A 2003}) A solid box is $15$ cm by $10$ cm by $8$ cm. A new solid is formed by removing a cube $3$ cm on a side from each corner of this box. What percent of the original volume is removed? 


 $\mathrm{(A) \ } 4.5\qquad \mathrm{(B) \ } 9\qquad \mathrm{(C) \ } 12\qquad \mathrm{(D) \ } 18\qquad \mathrm{(E) \ } 24$

 \textbf{Answer: }D



\textbf{Solution}

The volume of the original box is $15\cdot10\cdot8=1200.$ 

 The volume of each cube that is removed is $3\cdot3\cdot3=27.$

 Since there are $8$ corners on the box, $8$ cubes are removed. 

 So the total volume removed is $8\cdot27=216$. 

 Therefore, the desired percentage is $\frac{216}{1200}\cdot100 = \boxed{\mathrm{(D)}\ 18\%}.$



	\item (\textit{AMC 10 10A 2003}) It takes Mary $30$ minutes to walk uphill $1$ km from her home to school, but it takes her only $10$ minutes to walk from school to her home along the same route. What is her average speed, in km/hr, for the round trip? 


 $\mathrm{(A) \ } 3\qquad \mathrm{(B) \ } 3.125\qquad \mathrm{(C) \ } 3.5\qquad \mathrm{(D) \ } 4\qquad \mathrm{(E) \ } 4.5$

 \textbf{Answer: }A



\textbf{Solution 1}

Since she walked $1$ km to school and $1$ km back home, her total distance is $1+1=2$ km. 

 Since she spent $30$ minutes walking to school and $10$ minutes walking back home, her total time is $30+10=40$ minutes = $\frac{40}{60}=\frac{2}{3}$ hours. 

 Therefore her average speed in km/hr is $\frac{2}{\left(\frac{2}{3}\right)}=\boxed{\mathrm{(A)}\ 3}$.

 



\textbf{Solution 2}

The average speed of two speeds that travel the same distance is the \href{https://artofproblemsolving.com/wiki/index.php?title=Harmonic\_mean}{harmonic mean} of the speeds, or $\dfrac{2}{\left(\dfrac{1}{x}+\dfrac{1}{y}\right)}=\dfrac{2xy}{x+y}$ (for speeds $x$ and $y$). Mary's speed going to school is $2\,\text{km/hr}$, and her speed coming back is $6\,\text{km/hr}$. Plugging the numbers in, we get that the average speed is $\dfrac{2\times 6\times 2}{6+2}=\dfrac{24}{8}=\boxed{\mathrm{(A)}\ 3}$.

 


	\item (\textit{AMC 10 10A 2003}) Let $d$ and $e$ denote the solutions of $2x^{2}+3x-5=0$. What is the value of $(d-1)(e-1)$? 


 $\mathrm{(A) \ } -\frac{5}{2}\qquad \mathrm{(B) \ } 0\qquad \mathrm{(C) \ } 3\qquad \mathrm{(D) \ } 5\qquad \mathrm{(E) \ } 6$

 \textbf{Answer: }B



\textbf{Solution 1}

Using factoring: 

 $2x^{2}+3x-5=0$
$(2x+5)(x-1)=0$
$x = -\frac{5}{2}$ or $x=1$ 

 So $d$ and $e$ are $-\frac{5}{2}$ and $1$. 

 Therefore the answer is $\left(-\frac{5}{2}-1\right)(1-1)=\left(-\frac{7}{2}\right)(0)=\boxed{\mathrm{(B)}\ 0}$




\textbf{Solution 2}

We can use the sum and product of a quadratic (a.k.a Vieta):

 $(d-1)(e-1)=de-(d+e)+1 \implies\text{product}-\text{sum}+1 \implies \dfrac{c}{a}-\left(-\dfrac{b}{a}\right)+1 \implies \dfrac{b+c}{a}+1= \dfrac{5}{-5}+1=\boxed{\mathrm{(B)}\ 0}$




\textbf{Solution 3}

By inspection, we quickly note that $x=1$ is a solution to the equation, therefore the answer is

 $(d-1)(e-1)=(1-1)(e-1)=\boxed{\mathrm{(B)}\ 0}$




\textbf{Solution 4}

The form $(d-1)(e-1)$ resembles the factored form for the quadratic, namely $(x-d)(x-e)$ based on the information given. Note putting in 1 for $x$ in the that quadratic immediately yields the desired expression. Thus, $2(1)^2+3(1)-5 = \boxed{\mathrm{(B)}\ 0}$



	\item (\textit{AMC 10 10A 2003}) Define $x \heartsuit y$ to be $|x-y|$ for all real numbers $x$ and $y$. Which of the following statements is not true? 


 $\mathrm{(A) \ } x \heartsuit y = y \heartsuit x$ for all $x$ and $y$
$\mathrm{(B) \ } 2(x \heartsuit y) = (2x) \heartsuit (2y)$ for all $x$ and $y$
$\mathrm{(C) \ } x \heartsuit 0 = x$ for all $x$
$\mathrm{(D) \ } x \heartsuit x = 0$ for all $x$
$\mathrm{(E) \ } x \heartsuit y > 0$ if $x \neq y$

 \textbf{Answer: }C



\textbf{Solution}

We start by looking at the answers. Examining statement C, we notice:

 $x \heartsuit 0 = |x-0| = |x|$
$|x| \neq x$ when $x<0$, but statement C says that it does for all $x$. 

 Therefore the statement that is not true is $\boxed{\mathrm{(C)}\ x\heartsuit 0=x\ \text{for all}\ x}$



	\item (\textit{AMC 10 10A 2003}) How many non-congruent triangles with perimeter $7$ have integer side lengths? 


 $\mathrm{(A) \ } 1\qquad \mathrm{(B) \ } 2\qquad \mathrm{(C) \ } 3\qquad \mathrm{(D) \ } 4\qquad \mathrm{(E) \ } 5$

 \textbf{Answer: }B



\textbf{Solution}

By the \href{https://artofproblemsolving.com/wiki/index.php?title=Triangle\_inequality}{triangle inequality}, no side may have a length greater than the semiperimeter, which is $\frac{1}{2}\cdot7=3.5$. 

 Since all sides must be integers, the largest possible length of a side is $3$. Therefore, all such triangles must have all sides of length $1$, $2$, or $3$. Since $2+2+2=6<7$, at least one side must have a length of $3$. Thus, the remaining two sides have a combined length of $7-3=4$. So, the remaining sides must be either $3$ and $1$ or $2$ and $2$. Therefore, the number of triangles is $\boxed{\mathrm{(B)}\ 2}$.

 


	\item (\textit{AMC 10 10A 2003}) What is the probability that a randomly drawn positive factor of $60$ is less than $7$?


 $\mathrm{(A) \ } \frac{1}{10}\qquad \mathrm{(B) \ } \frac{1}{6}\qquad \mathrm{(C) \ } \frac{1}{4}\qquad \mathrm{(D) \ } \frac{1}{3}\qquad \mathrm{(E) \ } \frac{1}{2}$

 \textbf{Answer: }E



\textbf{Solution 1}

For any positive integer $n$ which is not a perfect square, exactly half of its positive factors will be less than $\sqrt{n}$, since each such factor can be paired with one that is larger than $\sqrt{n}$. (By contrast, if $n$ is a perfect square, one of its factors will be exactly $\sqrt{n}$, which would therefore have to be paired with itself.)

 Since $60$ is indeed not a perfect square, it follows that half of its positive factors are less than $\sqrt{60} \approx 7.746$. This estimate clearly shows that there are not even any integers, let alone factors of $60$, between $7$ and $\sqrt{60}$. Accordingly, exactly half of the positive factors of $60$ are in fact less than $7$, so the answer is precisely $\boxed{\mathrm{(E)}\ \frac{1}{2}}$.

 



\textbf{Solution 2}

Testing all positive integers less than $7$, we find that $1$, $2$, $3$, $4$, $5$, and $6$ all divide $60$. The prime factorization of $60$ is $2^2 \cdot 3 \cdot 5$, so using the standard \href{https://artofproblemsolving.com/wiki/index.php?title=Divisor\_function}{formula for the number of divisors}, the total number of divisors of $60$ is $3 \cdot 2 \cdot 2 = 12$. Therefore, the required probability is $\frac{6}{12} = \boxed{\mathrm{(E)}\ \frac{1}{2}}$.

 



\textbf{Solution 3 (brute force)}

Though this is not recommended for reasons of time, one can simply write out all the factors of $60$, eventually finding that 
 \[60 = 1 \cdot 60 = 2 \cdot 30 = 3 \cdot 20 = 4 \cdot 15 = 5 \cdot 12 = 6 \cdot 10.\] Hence $60$ has $12$ factors, of which $6$ are less than $7$ (namely, $1$, $2$, $3$, $4$, $5$, and $6$), so the answer is $\frac{6}{12} = \boxed{\mathrm{(E)}\ \frac{1}{2}}$.

 


	\item (\textit{AMC 10 10A 2003}) Simplify 


 $\sqrt[3]{x\sqrt[3]{x\sqrt[3]{x\sqrt{x}}}}$.


 $\mathrm{(A) \ } \sqrt{x}\qquad \mathrm{(B) \ } \sqrt[3]{x^{2}}\qquad \mathrm{(C) \ } \sqrt[27]{x^{2}}\qquad \mathrm{(D) \ } \sqrt[54]{x}\qquad \mathrm{(E) \ } \sqrt[81]{x^{80}}$

 \textbf{Answer: }A



\textbf{Solution 1}

$\sqrt[3]{x\sqrt{x}}=\sqrt[3]{(\sqrt{x})^{2}\cdot\sqrt{x}}=\sqrt[3]{(\sqrt{x})^{3}}=\sqrt{x}$. 

 Therefore: 

 $\sqrt[3]{x\sqrt[3]{x\sqrt[3]{x\sqrt{x}}}}=\sqrt[3]{x\sqrt[3]{x\sqrt{x}}}=\sqrt[3]{x\sqrt{x}}= \sqrt{x} \Rightarrow \boxed{\mathrm{(A)}\ \sqrt{x}}$




\textbf{Solution 2}

We know that $\sqrt[3]{x\sqrt{x}} = \sqrt[3]{x\cdot(x^\frac{1}{2})} = \sqrt[3]{x^\frac{3}{2}} = x^\frac{1}{2}$ We plug this into $\sqrt[3]{x\sqrt[3]{x\sqrt[3]{x\sqrt{x}}}}$ and get $\sqrt[3]{x\sqrt[3]{x\cdot(x^\frac{1}{2})}}$. Again, we substitute and get $\sqrt[3]{x\cdot(x^\frac{1}{2})}$. We substitute one more time and get $x^\frac{1}{2} = \boxed{(A) \sqrt{x}}$.

 



\textbf{Solution 3 (Lazy, but kind of intuitive)}

WLOG, plug in some random square number like $9$ or $4$, and the output you get after simplifying the expression will always be the square root of the number, so the answer is just $\boxed{(A) \sqrt{x}}$.

 



\textbf{Solution 3 (Substitution)}

Substitute $y=\sqrt{x}$. The expression then turns into $\sqrt[3]{y^2\sqrt[3]{y^2\sqrt[3]{y^2\sqrt{y^2}}}}$. Notice that for each cube root we get $\sqrt[3]{y^3}$ or $y$, so our final answer is simply $y$, or equivalently $\boxed{(A) \sqrt{x}}$.

 


	\item (\textit{AMC 10 10A 2003}) The polygon enclosed by the solid lines in the figure consists of 4 congruent squares joined edge-to-edge. One more congruent square is attached to an edge at one of the nine positions indicated. How many of the nine resulting polygons can be folded to form a cube with one face missing? 


 
                    \begin{figure}[H]
                        \centering
                        \includegraphics[width=0.5\textwidth]{Problem/AMC_10_10A_2003_10_Problem_Figure_1.png}
                    \end{figure}
                    
$\mathrm{(A) \ } 2\qquad \mathrm{(B) \ } 3\qquad \mathrm{(C) \ } 4\qquad \mathrm{(D) \ } 5\qquad \mathrm{(E) \ } 6$

 \textbf{Answer: }E



\textbf{Solution 1}


                    \begin{figure}[H]
                        \centering
                        \includegraphics[width=0.5\textwidth]{Solution/AMC_10_10A_2003_10_Solution_Figure_1.gif}
                    \end{figure}
                    

 Let the squares be labeled $A$, $B$, $C$, and $D$.

 When the polygon is folded, the "right" edge of square $A$ becomes adjacent to the "bottom edge" of square $C$, and the "bottom" edge of square $A$ becomes adjacent to the "bottom" edge of square $D$. 

 So, any "new" square that is attatched to those edges will prevent the polygon from becoming a cube with one face missing. 

 Therefore, squares $1$, $2$, and $3$ will prevent the polygon from becoming a cube with one face missing.

 Squares  $4$, $5$, $6$, $7$, $8$, and $9$ will allow the polygon to become a cube with one face missing when folded. 

 Thus the answer is $\boxed{\mathrm{(E)}\ 6}$.

 



\textbf{Solution 2}

Another way to think of it is that a cube missing one face has $5$ of its $6$ faces.  Since the shape has $4$ faces already, we need another face.  The only way to add another face is if the added square does not overlap any of the others.  $1$,$2$, and $3$ overlap, while squares $4$ to $9$ do not. The answer is $\boxed{\mathrm{(E)}\ 6}$




\textbf{Solution 3}

If you're good at visualizing, you can imagine each box and fold up the shape into a 3D shape. This solution is only recommended if you are either in a hurry or extremely skilled at visualizing. We find out that $4,5,6,7,8$ and $9$ work. Therefore, the answer is $\boxed{\mathrm{(E)}\ 6}$. ~Sophia866

 


	\item (\textit{AMC 10 10A 2003}) The sum of the two 5-digit numbers $AMC10$ and $AMC12$ is $123422$. What is $A+M+C$? 


 $\mathrm{(A) \ } 10\qquad \mathrm{(B) \ } 11\qquad \mathrm{(C) \ } 12\qquad \mathrm{(D) \ } 13\qquad \mathrm{(E) \ } 14$

 \textbf{Answer: }E



\textbf{Solution}

$AMC10+AMC12=123422$
$AMC00+AMC00=123400$
$AMC+AMC=1234$
$2\cdot AMC=1234$
$AMC=\frac{1234}{2}=617$

 Since $A$, $M$, and $C$ are digits, $A=6$, $M=1$, $C=7$. 

 Therefore, $A+M+C = 6+1+7 = \boxed{\mathrm{(E)}\ 14}$.

 



\textbf{Solution 2}

We know that $AMC12$ is $2$ more than $AMC10$. We set up $AMC10=x$ and $AMC12=x+2$. We have $x+x+2=123422$. Solving for $x$, we get $x=61710$. Therefore, the sum $A+M+C= 6+1+7 = \boxed{\mathrm{(E)}\ 14}$.

 



\textbf{Solution 3}

Consider the place values of the digits of $AMC10$ and $AMC12$. 

 When we add $AMC10$ and $AMC12$, $C + C$ must result in a units digit of $4$, meaning $C$ is either $2$ or $7$. Since $M$ is odd, this means a ten was carried over to the next place value from $C + C$, and thus $C = 7$ (as $7 + 7 = 14$ and the ten is carried over). Now, we know 3 is the units digit of $M + M + 1$, so $M$ is either $1$ or $6$. Again, we must look at the digit before $M$, or $A$. $A$ is even, so $M$ must be less than $5$, or else the ten would be carried over. Ergo, $M$ is $1$. Nothing is carried over, so we have $A + A = 12$, and $A = 6$. Therefore, the sum of $A$, $M$, and $C$ is $6 + 1 + 7 = \boxed{\mathrm{(E)}\ 14}$.

 



\textbf{Solution 4}

Intuitively, adding $AMC10$ and $AMC12$ is the same as adding $AMC11$ to itself. We can divide $123422$ by $2$ to get $61711$. So, $A$ = $6$, $M$ = 1, and $C$ = $7$. Adding all those up yields $6 + 1 + 7 = \boxed{\mathrm{(E)}\ 14}$ ~sdmd

 


	\item (\textit{AMC 10 10A 2003}) A point $(x,y)$ is randomly picked from inside the rectangle with  vertices $(0,0)$, $(4,0)$, $(4,1)$, and $(0,1)$. What is the probability that $x<y$? 


 $\mathrm{(A) \ } \frac{1}{8}\qquad \mathrm{(B) \ } \frac{1}{4}\qquad \mathrm{(C) \ } \frac{3}{8}\qquad \mathrm{(D) \ } \frac{1}{2}\qquad \mathrm{(E) \ } \frac{3}{4}$

 \textbf{Answer: }A



\textbf{Solution}

The rectangle has a width of $4$ and a height of $1$. 

 The area of this rectangle is $4\cdot1=4$. 

 The line $x=y$ intersects the rectangle at $(0,0)$ and $(1,1)$. 

 The area which $x<y$ is the right isosceles triangle with side length $1$ that has vertices at $(0,0)$, $(1,1)$, and $(0,1)$.  

 The area of this triangle is $\frac{1}{2}\cdot1^{2}=\frac{1}{2}$ 

 Therefore, the probability that $x<y$ is $\frac{\frac{1}{2}}{4}=\frac{1}{8} \Rightarrow \boxed{\mathrm{(A)}\ \frac{1}{8}}$



	\item (\textit{AMC 10 10A 2003}) The sum of three numbers is $20$. The first is four times the sum of the other two. The second is seven times the third. What is the product of all three? 


 $\mathrm{(A) \ } 28\qquad \mathrm{(B) \ } 40\qquad \mathrm{(C) \ } 100\qquad \mathrm{(D) \ } 400\qquad \mathrm{(E) \ } 800$

 \textbf{Answer: }A



\textbf{Solution 1}

Let the numbers be $x$, $y$, and $z$ in that order. The given tells us that 

 
 \begin{eqnarray*}y&=&7z\\ x&=&4(y+z)=4(7z+z)=4(8z)=32z\\ x+y+z&=&32z+7z+z=40z=20\\ z&=&\frac{20}{40}=\frac{1}{2}\\ y&=&7z=7\cdot\frac{1}{2}=\frac{7}{2}\\ x&=&32z=32\cdot\frac{1}{2}=16 \end{eqnarray*}

 Therefore, the product of all three numbers is $xyz=16\cdot\frac{7}{2}\cdot\frac{1}{2}=28 \Rightarrow \boxed{\mathrm{(A)}\ 28}$.

 



\textbf{Solution 2}

Alternatively, we can set up the system in \href{https://artofproblemsolving.com/wiki/index.php?title=Equation}{equation} form:

 
 \begin{eqnarray*}1x+1y+1z&=&20\\ 1x-4y-4z&=&0\\ 0x+1y-7z&=&0\\ \end{eqnarray*}

 Or, in \href{https://artofproblemsolving.com/wiki/index.php?title=Matrix}{matrix} form$\begin{bmatrix} 1 & 1 & 1 \\ 1 & -4 & -4 \\ 0 & 1 & -7 \end{bmatrix} \begin{bmatrix} x \\ y \\ z \\ \end{bmatrix} =\begin{bmatrix} 20 \\ 0 \\ 0 \\ \end{bmatrix}$

 To solve this matrix equation, we can rearrange it thus:

 $\begin{bmatrix} x \\ y \\ z \\ \end{bmatrix} = \begin{bmatrix} 1 & 1 & 1 \\ 1 & -4 & -4 \\ 0 & 1 & -7 \end{bmatrix} ^{-1} \begin{bmatrix} 20 \\ 0 \\ 0 \\ \end{bmatrix}$

 Solving this matrix equation by using \href{https://artofproblemsolving.com/wiki/index.php?title=Inverse\_matrices&action=edit&redlink=1}{inverse matrices} and \href{https://artofproblemsolving.com/wiki/index.php?title=Matrix\_multiplication&action=edit&redlink=1}{matrix multiplication} yields

 $\begin{bmatrix} x \\ y \\ z \\ \end{bmatrix} = \begin{bmatrix} \frac{1}{2} \\ \frac{7}{2} \\ 16 \\ \end{bmatrix}$

 Which means that $x = \frac{1}{2}$, $y = \frac{7}{2}$, and $z = 16$. Therefore, $xyz = \frac{1}{2}\cdot\frac{7}{2}\cdot16 = 28 \Rightarrow \boxed{\mathrm{(A)}\ 28}$




\textbf{Solution 3}

Let's denote the 3rd number as $x$, the 2nd as $7x$, and the 1st as $4(7x+x)$ according to the information given in the problem. We know that all three numbers add up to $20$, so $4(8x)+7x+x = 20$. Solving the equation we get $x = \frac{1}{2}$. Then substitute this value of x to solve for the other two numbers. Lastly, we obtain $28$ as the product of all three numbers.

 


	\item (\textit{AMC 10 10A 2003}) Let $n$ be the largest integer that is the product of exactly 3 distinct prime numbers $d$, $e$, and $10d+e$, where $d$ and $e$ are single digits. What is the sum of the digits of $n$?


 $\mathrm{(A) \ } 12\qquad \mathrm{(B) \ } 15\qquad \mathrm{(C) \ } 18\qquad \mathrm{(D) \ } 21\qquad \mathrm{(E) \ } 24$

 \textbf{Answer: }A



\textbf{Solution 1}

Since we want $n$ to be as large as possible, we would like $d$ in $10d+e$ to be as large as possible. So, $d=7,$ the greatest single-digit prime. Then, $e$ cannot be $5$ because $10(7)+5 = 75,$ which is not prime. So $e = 3$. Therefore, $d \cdot e \cdot (10d+e) = 7 \cdot 3 \cdot 73 = 1533$.So, the sum of the digits of $n$ is $1+5+3+3=12 \Rightarrow \boxed{\mathrm{(A)}\ 12}$ ~ MathGenius\_ (Edited by Sophia866)

 



\textbf{Solution 2}

Since $d$ is a single digit prime number, the set of possible values of $d$ is $\{2,3,5,7\}$. 

 Since $e$ is a single digit prime number and is the units digit of the prime number $10d+e$, the set of possible values of $e$ is $\{3,7\}$. 

 Using these values for $d$ and $e$, the set of possible values of $10d+e$ is $\{23,27,33,37,53,57,73,77\}$

 Out of this set, the prime values are $\{23,37,53,73\}$

 Therefore the possible values of $n$ are: 

 $2\cdot3\cdot23=138$
$3\cdot7\cdot37=777$
$5\cdot3\cdot53=795$
$7\cdot3\cdot73=1533$ 

 The largest possible value of $n$ is $1533$. 

 So, the sum of the digits of $n$ is $1+5+3+3=12 \Rightarrow \boxed{\mathrm{(A)}\ 12}$



	\item (\textit{AMC 10 10A 2003}) What is the probability that an integer in the set $\{1,2,3,...,100\}$ is divisible by $2$ and not divisible by $3$? 


 $\mathrm{(A) \ } \frac{1}{6}\qquad \mathrm{(B) \ }  \frac{33}{100}\qquad \mathrm{(C) \ }  \frac{17}{50}\qquad \mathrm{(D) \ }  \frac{1}{2}\qquad \mathrm{(E) \ }  \frac{18}{25}$

 \textbf{Answer: }C



\textbf{Solution}

There are $100$ integers in the set. 

 Since every $2^{\text{nd}}$ integer is divisible by $2$, there are $\left\lfloor\frac{100}{2}\right\rfloor = 50$ integers divisible by $2$ in the set. 

 To be divisible by both $2$ and $3$, a number must be divisible by $\operatorname{lcm}(2,3)=6$.

 Since every $6^{\text{th}}$ integer is divisible by $6$, there are $\left\lfloor\frac{100}{6}\right\rfloor = 16$ integers divisible by both $2$ and $3$ in the set. 

 So there are $50-16=34$ integers in this set that are divisible by $2$ and not divisible by $3$. 

 Therefore, the desired probability is $\frac{34}{100}=\frac{17}{50}\Rightarrow\boxed{\mathrm{(C)}\ \frac{17}{50}}$.

 


	\item (\textit{AMC 10 10A 2003}) What is the units digit of $13^{2003}$?


 $\mathrm{(A) \ } 1\qquad \mathrm{(B) \ } 3\qquad \mathrm{(C) \ } 7\qquad \mathrm{(D) \ } 8\qquad \mathrm{(E) \ } 9$

 \textbf{Answer: }C



\textbf{Solution 1}

$13^{2003}\equiv 3^{2003}\pmod{10}$

 Since $3^4=81\equiv1\pmod{10}$: 

 $3^{2003}=(3^{4})^{500}\cdot3^{3}\equiv1^{500}\cdot27\equiv7\pmod{10}$

 Therefore, the units digit is $7 \Rightarrow\boxed{\mathrm{(C)}\ 7}$




\textbf{Solution 2 (Patterns)}

Since we are looking for the units digit of $13^{2003}$, we only have to focus on the units digit of the base (13) as none of the other digits of the base affect the units digit of the resulting value.

 By testing the first few values or through previous knowledge, you might see that the units digit of exponents with base 3 follow this pattern: 
 \[3^1=3\]
 \[3^2=9\] 
 \[3^3=27\]
 \[3^4=81,\]giving us the rotation $3-9-7-1.$

 As this cycle resets every time the index increases by 4, we know that this cycle ends on 2000, and starts once again on 2001. As our expression is raised to the power of 2003, we know that the units digit of our expression must end with the third term of our pattern: $7$.

 Therefore, the units digit of our expression is $7 \Rightarrow\boxed{\mathrm{(C)}\ 7}$



	\item (\textit{AMC 10 10A 2003}) The number of inches in the perimeter of an equilateral triangle equals the number of square inches in the area of  its circumscribed circle. What is the radius, in inches, of the circle? 


 $\mathrm{(A) \ } \frac{3\sqrt{2}}{\pi}\qquad \mathrm{(B) \ }  \frac{3\sqrt{3}}{\pi}\qquad \mathrm{(C) \ } \sqrt{3}\qquad \mathrm{(D) \ } \frac{6}{\pi}\qquad \mathrm{(E) \ } \sqrt{3}\pi$

 \textbf{Answer: }B



\textbf{Solution 1}

Let $s$ be the length of a side of the equilateral triangle and let $r$ be the radius of the circle. 

 In a circle with a radius $r$, the side of an inscribed equilateral triangle is $r\sqrt{3}$. 

 So $s=r\sqrt{3}$. 

 The perimeter of the triangle is $3s=3r\sqrt{3}$

 The area of the circle is $\pi r^{2}$

 So: $\pi r^{2} = 3r\sqrt{3}$
$\pi r=3\sqrt{3}$
$r=\frac{3\sqrt{3}}{\pi} \Rightarrow\boxed{\mathrm{(B)}\ \frac{3\sqrt{3}}{\pi}}$




\textbf{Solution 2}

As in Solution $1$, let $s$ be the side length of the equilateral triangle and let $r$ be the radius of the circumcircle. Remembering that the formula for radius of a circumcircle is $\frac{abc}{4A}$, where $a$, $b$, and $c$ are side lengths of the triangle, and $A$ is its area, we define each of the quantities in the formula. Since $a = b = c = s$ and, by the formula for the area of an equilateral triangle, $A = \frac{s^2\sqrt{3}}{4}$, the circumradius is $\frac{s^3}{4\frac{s^2\sqrt{3}}{4}} = \frac{s}{\sqrt{3}}$. The area of the circumcircle is then $\pi(\frac{s}{\sqrt{3}})^2 = \frac{\pi s^2}{3}$. Because this is equal to the perimeter of the triangle, we equate this to $3s$. Upon solving, we note that the nonzero value of $s$ is then $\frac{9}{\pi}$. We can now plug this back into the circumradius formula we found earlier to get that $r = \frac{\frac{9}{\pi}}{\sqrt{3}} = \boxed{\textbf{(B)} \frac{3 \sqrt{3}}{\pi}}$.

 


	\item (\textit{AMC 10 10A 2003}) What is the sum of the reciprocals of the roots of the equation 


 $\frac{2003}{2004}x+1+\frac{1}{x}=0$?


 $\mathrm{(A) \ } -\frac{2004}{2003}\qquad \mathrm{(B) \ } -1\qquad \mathrm{(C) \ } \frac{2003}{2004}\qquad \mathrm{(D) \ } 1\qquad \mathrm{(E) \ } \frac{2004}{2003}$

 \textbf{Answer: }B



\textbf{Solution 1}

Multiplying both sides by $x$: 

 $\frac{2003}{2004}x^{2}+x+1=0$

 Let the roots be $a$ and $b$. 

 The problem is asking for $\frac{1}{a}+\frac{1}{b}= \frac{a+b}{ab}$

 By \href{https://artofproblemsolving.com/wiki/index.php?title=Vieta\%27s\_Formulas}{Vieta's Formulas}: 

 $a+b=(-1)^{1}\frac{1}{\frac{2003}{2004}}=-\frac{2004}{2003}$
$ab=(-1)^{2}\frac{1}{\frac{2003}{2004}}=\frac{2004}{2003}$

 So the answer is $\frac{a+b}{ab}=\frac{-\frac{2004}{2003}}{\frac{2004}{2003}}=-1 \Rightarrow\boxed{\mathrm{(B)}\ -1}$.

 



\textbf{Solution 2}

Dividing both sides by $x$, 

 $\frac{2003}{2004}+\frac 1x+\frac{1}{x^2}=0$, 

 we see by \href{https://artofproblemsolving.com/wiki/index.php?title=Vieta\%27s\_Formulas}{Vieta's Formulas} that the sum of the roots is $-1  \Rightarrow\boxed{\mathrm{(B)}\ -1}$.

 



\textbf{Solution 3(Alternate Approach to Solution 1)}

Re-stating and multiplying $\frac{1}{x}$ by $\frac{2004}{2004}$, we have $\frac{2003x}{2004} + \frac{2004}{2004x} + 1 = 0$.

 Simplifying, we have $\frac{2003}{2004}+\frac{\frac{2004}{x}}{2004}=-1.$

 Putting this together, we have $\frac{2003+\frac{2004}{x}}{2004}=-1.$

 Simplifying, we have $2004x + \frac{2004}{x} = -2004.$

 Multiplying both sides by $x,$ we have $2003x^2+2004=-2004x.$

 Moving $-2004x$ to the left side, we have $2003x^2 + 2004x + 2004 = 0.$

 Let's go back and see what we want.

 We want the sum of the reciprocals of the roots, in which if the roots were to be called $r_1$ and $r_2,$ we would want $\frac{1}{r_1}+\frac{1}{r_2}=\frac{r_1 + r_2}{r_1 \cdot r_2}$.

 We can use \href{https://artofproblemsolving.com/wiki/index.php?title=Vieta\%27s\_Formulas}{Vieta's Formulas} to solve the sum of the roots and the product of the roots.

 We have that the sum of the two roots is $\frac{-b}{a}$ where the quadratic is $ax^2+bx+c=0.$

 In this case, $a=2003$ and $b=2004$.

 Therefore, $r_1+r_2=-\frac{2004}{2003}.$

 We have that the product of the two roots is $\frac{c}{a}$ where the quadratic is $ax^2+bx+c=0.$

 In this case, $c=2004$ and $a=2003.$

 Therefore, $r_1 \cdot r_2 = \frac{2004}{2003}.$

 Now that we have the sum and product of the roots, we can substitute back into the expression of what we want.

 We have $\frac{-\frac{2004}{2003}}{\frac{2004}{2003}}=-\frac{2004}{2003} \cdot \frac{2003}{2004} = \boxed{\mathrm{(B)}\ -1}.$



	\item (\textit{AMC 10 10A 2003}) A semicircle of diameter $1$ sits at the top of a semicircle of diameter $2$, as shown. The shaded area inside the smaller semicircle and outside the larger semicircle is called a \textit{lune}. Determine the area of this lune. 


 
                    \begin{figure}[H]
                        \centering
                        \includegraphics[width=0.5\textwidth]{Problem/AMC_10_10A_2003_19_Problem_Figure_1.png}
                    \end{figure}
                    
$\mathrm{(A) \ } \frac{1}{6}\pi-\frac{\sqrt{3}}{4}\qquad \mathrm{(B) \ } \frac{\sqrt{3}}{4}-\frac{1}{12}\pi\qquad \mathrm{(C) \ } \frac{\sqrt{3}}{4}-\frac{1}{24}\pi\qquad \mathrm{(D) \ } \frac{\sqrt{3}}{4}+\frac{1}{24}\pi\qquad \mathrm{(E) \ } \frac{\sqrt{3}}{4}+\frac{1}{12}\pi$

 \textbf{Answer: }C



\textbf{Solution 1}


                    \begin{figure}[H]
                        \centering
                        \includegraphics[width=0.5\textwidth]{Solution/AMC_10_10A_2003_19_Solution_Figure_1.png}
                    \end{figure}
                    

 
The shaded area $[A]$ is equal to the area of the smaller semicircle $[A+B]$ minus the area of a \href{https://artofproblemsolving.com/wiki/index.php?title=Sector}{sector} of the larger circle $[B+C]$ plus the area of a \href{https://artofproblemsolving.com/wiki/index.php?title=Triangle}{triangle} formed by two \href{https://artofproblemsolving.com/wiki/index.php?title=Radius}{ radii} of the larger semicircle and the diameter of the smaller semicircle $[C]$.

 The area of the smaller semicircle is $[A+B] = \frac{1}{2}\pi\cdot\left(\frac{1}{2}\right)^{2}=\frac{1}{8}\pi$. 

 Since the radius of the larger semicircle is equal to the diameter of the smaller half circle, the triangle is an equilateral triangle and the sector measures $60^\circ$. 

 The area of the $60^\circ$ sector of the larger semicircle is $[B+C] = \frac{60}{360}\pi\cdot\left(\frac{2}{2}\right)^{2}=\frac{1}{6}\pi$. 

 The area of the triangle is $[C] = \frac{1^{2}\sqrt{3}}{4}=\frac{\sqrt{3}}{4}$.

 So the shaded area is $[A] = [A+B]-[B+C]+[C] = \left(\frac{1}{8}\pi\right)-\left(\frac{1}{6}\pi\right)+\left(\frac{\sqrt{3}}{4}\right)=\boxed{\mathrm{(C)}\ \frac{\sqrt{3}}{4}-\frac{1}{24}\pi}$. We have thus solved the problem.

 


	\item (\textit{AMC 10 10A 2003}) A base-10 three digit number $n$ is selected at random. Which of the following is closest to the probability that the base-9 representation and the base-11 representation of $n$ are both three-digit numerals? 


 $\mathrm{(A) \ } 0.3\qquad \mathrm{(B) \ } 0.4\qquad \mathrm{(C) \ } 0.5\qquad \mathrm{(D) \ } 0.6\qquad \mathrm{(E) \ } 0.7$

 \textbf{Answer: }E



\textbf{Solution}

To be a three digit number in base-10:  

 $10^{2} \leq n \leq 10^{3}-1$
$100 \leq n \leq 999$ 

 Thus there are $900$ three-digit numbers in base-10

 To be a three-digit number in base-9:  

 $9^{2} \leq n \leq 9^{3}-1$
$81 \leq n \leq 728$

 To be a three-digit number in base-11:  

 $11^{2} \leq n \leq 11^{3}-1$
$121 \leq n \leq 1330$

 So, $121 \leq n \leq 728$ 

 Thus, there are $608$ base-10 three-digit numbers that are three digit numbers in base-9 and base-11. 

 Therefore the desired probability is $\frac{608}{900}\approx 0.7 \Rightarrow\boxed{\mathrm{(E)}\ 0.7}$.

 


	\item (\textit{AMC 10 10A 2003}) Pat is to select six cookies from a tray containing only chocolate chip, oatmeal, and peanut butter cookies. There are at least six of each of these three kinds of cookies on the tray. How many different assortments of six cookies can be selected? 


 $\mathrm{(A) \ } 22\qquad \mathrm{(B) \ } 25\qquad \mathrm{(C) \ } 27\qquad \mathrm{(D) \ } 28\qquad \mathrm{(E) \ } 729$

 \textbf{Answer: }D



\textbf{Solution 1 (Stars and Bars)}

It is given that it is possible to select at least 6 of each. Therefore, we can make a bijection to the number of ways to divide the six choices into three categories, since it is assumed that their order is unimportant. Using the ball and urns/sticks and stones/stars and bars formula, the number of ways to do this is $\binom{6+3-1}{3-1}=\binom{8}{2} = 28 \Rightarrow\boxed{\mathrm{(D)}\ 28}$~hi

 



\textbf{Solution 2 (Casework)}

Let the ordered triplet $(x,y,z)$ represent the assortment of $x$ chocolate chip cookies, $y$ oatmeal cookies, and $z$ peanut butter cookies. 

 Using \href{https://artofproblemsolving.com/wiki/index.php?title=Casework}{casework}: 

 Pat selects $0$ chocolate chip cookies: 

 Pat needs to select $6-0=6$ more cookies that are either oatmeal or peanut butter. 

 The assortments are: $\{(0,6,0); (0,5,1); (0,4,2); (0,3,3); (0,2,4); (0,1,5); (0,0,6)\} \rightarrow 7$ assortments. 

 Pat selects $1$ chocolate chip cookie: 

 Pat needs to select $6-1=5$ more cookies that are either oatmeal or peanut butter. 

 The assortments are: $\{(1,5,0); (1,4,1); (1,3,2); (1,2,3); (1,1,4); (1,0,5) \} \rightarrow 6$ assortments. 

 Pat selects $2$ chocolate chip cookies: 

 Pat needs to select $6-2=4$ more cookies that are either oatmeal or peanut butter. 

 The assortments are: $\{(2,4,0); (2,3,1); (2,2,2); (2,1,3); (2,0,4)\} \rightarrow 5$ assortments. 

 Pat selects $3$ chocolate chip cookies: 

 Pat needs to select $6-3=3$ more cookies that are either oatmeal or peanut butter. 

 The assortments are: $\{(3,3,0); (3,2,1); (3,1,2); (3,0,3)\} \rightarrow 4$ assortments. 

 Pat selects $4$ chocolate chip cookies: 

 Pat needs to select $6-4=2$ more cookies that are either oatmeal or peanut butter. 

 The assortments are: $\{(4,2,0); (4,1,1); (4,0,2)\} \rightarrow 3$ assortments. 

 Pat selects $5$ chocolate chip cookies: 

 Pat needs to select $6-5=1$ more cookies that are either oatmeal or peanut butter. 

 The assortments are: $\{(5,1,0); (5,0,1)\} \rightarrow 2$ assortments. 

 Pat selects $6$ chocolate chip cookies: 

 Pat needs to select $6-6=0$ more cookies that are either oatmeal or peanut butter. 

 The only assortment is: $\{(6,0,0)\} \rightarrow 1$ assortment. 

 The total number of assortments of cookies that can be collected is $7+6+5+4+3+2+1=28 \Rightarrow\boxed{\mathrm{(D)}\ 28}$

 There is a much faster way to do casework.

 Case 1: 1 type of cookie - 3 ways

 Case 2: 2 types of cookies - 3 ways to choose the 2 types of cookies, and 5 ways to choose how of each there are1 and 52 and 43 and 34 and 25 and 1

 $3*5=15$ cases for case 2

 Case 3: 3 types of cookies - quick examination shows us that the only ways to use all three cookies are the following:4, 1, 1: this gives us 3!/2!*1! = 3 ways3, 2, 1: this gives us 3! = 6 ways2, 2, 2: this gives us 1 waytotal of 10 cases for this case10+15+3= $28$ total

 


 



\textbf{Solution 3}

Suppose the six cookies to be chosen are the stars, as we attempt to implement stars and bars. We take two dividers, and place them between the cookies, such that the six cookies are split into 3 groups, where the groups are the number of chocolate chip, oatmeal, and peanut butter cookies, and each group can have 0. First, assume that the two dividers cannot go in between the same two cookies. By stars and bars, there are $\dbinom{7}{2}$ ways to make the groups.

 Finally, since the two dividers can be together, we must add those cases where the two dividers are in the same space between cookies. There are 7 spaces, and hence 7 cases.

 Our final answer is $21+7=\boxed{\mathrm{(D)}\ 28}$




\textbf{Solution 4(Alcumus Solution 1)}

The numbers of the three types of cookies must have a sum of six. Possible sets of whole numbers whose sum is six are 
 \[0,0,6;\ 0,1,5;\ 0,2,4;\ 0,3,3;\ 1,1,4;\ 1,2,3;\ \ \text{and}\ 2,2,2.\] Every ordering of each of these sets determines a different assortment of cookies. There are 3 orders for each of the sets 
 \[0,0,6;\ 0,3,3;\ \text{and}\ 1,1,4.\] There are 6 orders for each of the sets 
 \[0,1,5;\ 0,2,4;\ \text{and}\ 1,2,3.\] There is only one order for $2,2,2$. Therefore the total number of assortments of six cookies is $3\cdot 3 + 3\cdot 6 + 1 = 28$.

 



\textbf{Solution 5(Alcumus Solution 2)}

Also known as the \href{https://artofproblemsolving.com/wiki/index.php?title=Ball-and-urn}{ball-and-urn} technique.Construct eight slots, six to place the cookies in and two to divide the cookies by type. Let the number of chocolate chip cookies be the number of slots to the left of the first divider, the number of oatmeal cookies be the number of slots between the two dividers, and the number of peanut butter cookies be the number of slots to the right of the second divider. For example, $111\ | \ 11\ | \ 1$ represents three chocolate chip cookies, two oatmeal cookies, and one peanut butter cookie. There are $\binom{8}{2} = \boxed{28}$ ways to place the two dividers, so there are 28 ways to select the six cookies.

 \textbf{Cheap Solution 6} Notice that the answer will be able to be represented by a choose binomial $\binom{n}{k}$, where $n$ is a bit greater than six. Looking at the answer choices, we see $\boxed{28}$.

 



\textbf{Solution 7}

(NOTE: Variation of other solutions)

 Let $a, b,c$ be the number of cookies that we have for each type respectively. Hence:

 $a + b + c = 6.$

 Via stars and bars, the number of solutions is simply $\tbinom{8}{2} = 28.$




\textbf{Solution 8}

This is simply a stars and bars situation, where we distribute 6 items among 3 people. 

 A quick formula for stars and bars problems is $\tbinom{n+k-1}{k-1}$ where we use combinatrics to evaluate after.Using this, we get $\tbinom{8}{2} = 28.$.

 So, the answer is (D) 28.

 


	\item (\textit{AMC 10 10A 2003}) In rectangle $ABCD$, we have $AB=8$, $BC=9$, $H$ is on $BC$ with $BH=6$, $E$ is on $AD$ with $DE=4$, line $EC$ intersects line $AH$ at $G$, and $F$ is on line $AD$ with $GF \perp AF$. Find the length of $GF$. 


 
                    \begin{figure}[H]
                        \centering
                        \includegraphics[width=0.5\textwidth]{Problem/AMC_10_10A_2003_22_Problem_Figure_1.png}
                    \end{figure}
                    
$\mathrm{(A) \ } 16\qquad \mathrm{(B) \ } 20\qquad \mathrm{(C) \ } 24\qquad \mathrm{(D) \ } 28\qquad \mathrm{(E) \ } 30$

 \textbf{Answer: }B



\textbf{Solution 1}

$\angle GHC = \angle AHB$ (Vertical angles are equal).

 $\angle F = \angle B$ (Both are 90 degrees).

 $\angle BHA = \angle HAD$ (Alt. Interior Angles are congruent).

 Therefore $\triangle GFA$ and $\triangle ABH$ are similar.$\triangle GCH$ and $\triangle GEA$ are also similar.

 $DA$ is 9, therefore $EA$ must equal 5. Similarly, $CH$ must equal 3. 

 Because $GCH$ and $GEA$ are similar, the ratio of $CH\; =\; 3$ and $EA\; =\; 5$, must also hold true for $GH$ and $HA$. $\frac{GH}{GA} = \frac{3}{5}$, so $HA$ is $\frac{2}{5}$ of $GA$. By Pythagorean theorem, $(HA)^2\;  =\; (HB)^2\; +\; (BA)^2\;...\;HA=10$.

 $HA\: =\: 10 =\: \frac{2}{5}*(GA)$. 

 $GA\: =\: 25.$

 So $\frac{GA}{HA}\: =\: \frac{GF}{BA}$.

 $\frac{25}{10}\: =\: \frac{GF}{8}$.

 Therefore   $GF= \boxed{\mathrm{(B)}\ 20}$.

 



\textbf{Solution 2}

Since $ABCD$ is a rectangle, $CD=AB=8$. 

 Since $ABCD$ is a rectangle and $GF \perp AF$,  $\angle GFE = \angle CDE = \angle ABC = 90^\circ$. 

 Since $ABCD$ is a rectangle, $AD || BC$. 

 So, $AH$ is a \href{https://artofproblemsolving.com/wiki/index.php?title=Transversal}{transversal}, and $\angle GAF = \angle AHB$. 

 This is sufficient to prove that $GFE \approx CDE$ and $GFA \approx ABH$.

 Using ratios: 

 $\frac{GF}{FE}=\frac{CD}{DE}$
$\frac{GF}{FD+4}=\frac{8}{4}=2$
$GF=2 \cdot (FD+4)=2 \cdot FD+8$
$\frac{GF}{FA}=\frac{AB}{BH}$
$\frac{GF}{FD+9}=\frac{8}{6}=\frac{4}{3}$
$GF=\frac{4}{3} \cdot (FD+9)=\frac{4}{3} \cdot FD+12$

 Since $GF$ can't have 2 different lengths, both expressions for $GF$ must be equal. 

 $2 \cdot FD+8=\frac{4}{3} \cdot FD+12$
$\frac{2}{3} \cdot FD=4$
$FD=6$
$GF=2 \cdot FD+8=2\cdot6+8=\boxed{\mathrm{(B)}\ 20}$




\textbf{Solution 3 (fastest and easiest)}

We extend $BC$ such that it intersects $GF$ at $X$. Since $ABCD$ is a rectangle, it follows that $CD=8$, therefore, $XF=8$. Let $GX=y$. From the similarity of triangles $GCH$ and $GEA$, we have the ratio $3:5$ (as $CH=9-6=3$, and $EA=9-4=5$). $GX$ and $GF$ are the altitudes of $GCH$ and $GEA$, respectively. Thus, $y:y+8 = 3:5$, from which we have $y=12$, thus $GF=y+8=12+8=\boxed{\mathrm{(B)}\ 20}$




\textbf{Solution 4}

Since $GF\perp AF$ and $AF\perp CD,$ we have $GF\parallel CD\parallel AB.$ Thus, $\triangle CDE\sim GFE.$ Suppose $GF=x$ and $FD=y.$ Thus, we have $\dfrac{x}{8}=\dfrac{y+4}{4}.$ Additionally, now note that $\triangle GAF\sim AHB,$ which is pretty obvious from insight, but can be proven by AA with extending $BH$ to meet $GF.$ From this new pair of similar triangles, we have $\dfrac{x}{8}=\dfrac{y+9}{6}.$ Therefore, we have by combining those two equations, 
 \[\dfrac{y+9}{6}=\dfrac{y+4}{4}.\] Solving, we have $y=6,$ and therefore $x=\boxed{\mathrm{(B)}\ 20}$




\textbf{Solution 5}

Since there are only lines, you can resort to coordinate bashing. Let $FD=k$. Three lines, line $GF$, line $GE$, and line $GA$, intersect at $G$. Our goal is to find the y-coordinate of that intersection point. 

 Line $GF$ is $x=0$

 Line $GE$ passes through $(k+4, 0)$ and $(k, 8)$. Therefore the slope is $-2$ and the line is $y-0=-2(x-k-4)$ which is $y=-2x+2k+8$

 Line $GA$ passes through $(k+9, 0)$ and $(k+3, 8)$. Therefore the slope is $\frac{-4}{3}$ and the line is$y-0=-\frac{4}{3}(x-k-9)$ which simplifies to $y=-\frac{4}{3}x+\frac{4}{3}k+12$

 We solve the system of equations with these three lines. First we plug in $x=0$
$y=2k+8$
$y=\frac{4}{3}k+12$

 Next, we solve for k. $k=6$ Therefore $y=20$. The y-coordinate of this intersection point is indeed our answer. $\boxed{\mathrm{(B)}\ 20}$~superagh

 



\textbf{Solution 6 (simple coordinates)}

Let $A$ be the origin of our coordinate system. Now line $GA$ has equation $-\frac{4}{3}x$. We can use point-slope form to find the equation for line $GE$. First, we know that its slope is $-2$, and we know that it passes through $E=(-5,0)$, so line $GE$ has equation $-2(x+5)$. Solving for the intersection by letting $-\frac{4}{3}x=-2(x+5)$, we get $x=-15$. Plugging this into our equation for line $GA$ gives us $G=(-15,20)$, so $GF= \boxed{\mathrm{(B)}\ 20}$~chrisdiamond10

 



\textbf{Solution 7 (system of equations through angle similarity)}

First, using given information, we can find the values of some line segments in the figure. We find that $HA = 10$ (through Pythagorean Theorem), $CH = 3$, and $EA = 5$. Let Line $FD = x$ and let Line $FG = y$. We find that $\triangle FGE \sim \triangle CDE$ through some angle chasing (they both have a right angle, and they both share angle $\angle CED$. Using this information, we can write the equation $\frac{4}{8} = \frac{4+x}{y}$. Through simplifying this equation, we get that $y=2x+8$. Let point $I$ be the point on line $FG$ so that lines $CI$ and $FG$ are perpendicular, and we get that $GI = 2x$ and $FI =8$. Doing some more angle chasing, we can find that $\triangle GIH \sim \triangle GFA$, as they both share $\angle FGH$ and they both have a right angle.

 With this information, we can write the equation $\frac{x+3}{y-8} = \frac{x+9}{y}.$ Simplifying this equation we get the equation $-8x+6y-72 = 0$.Plugging in $y=2x+8$ for $6y$, we get $4x-24 = 0$, so $x=6$.Lastly, to find the value of y, which is the value of Line $FG$, our desired value, we plug in $6$ for $x$ in the equation $y=2x+8$, we get $2(6)+8$, which, finally, we get our $y$ value of $20$, so therefore, our answer is $GF= \boxed{\mathrm{(B)}\ 20}$




\textbf{Solution 8 Line and Slope}

Draw a coordinate plane with the y-axis centered on $CD$ and the x-axis centered on $AD$. From there, call the line passing through $CE$ $a$ and the line passing through $HA$ $b$. From there, you can find that the equations for these lines are $y=-2x+8$ and $y=-4/3 x+12$ respectively. We need to find the length of $GF$, so we are finding the y-value of the point $G$. Solving for this point, we get $-4/3 x+12=-2x+8$;$2/3x=-4$;$x=-6$. Now, plugging in the values, we find that $y=-2(-6)+8=20$, so our answer is $GF= \boxed{\mathrm{(B)}\ 20}$.

 


	\item (\textit{AMC 10 10A 2003}) A large equilateral triangle is constructed by using toothpicks to create rows of small equilateral triangles. For example, in the figure we have $3$ rows of small congruent equilateral triangles, with $5$ small triangles in the base row. How many toothpicks would be needed to construct a large equilateral triangle if the base row of the triangle consists of $2003$ small equilateral triangles? 


 
                    \begin{figure}[H]
                        \centering
                        \includegraphics[width=0.5\textwidth]{Problem/AMC_10_10A_2003_23_Problem_Figure_1.png}
                    \end{figure}
                    
$\mathrm{(A) \ } 1,004,004 \qquad \mathrm{(B) \ } 1,005,006 \qquad \mathrm{(C) \ } 1,507,509 \qquad \mathrm{(D) \ } 3,015,018 \qquad \mathrm{(E) \ } 6,021,018$

 \textbf{Answer: }C



\textbf{Solution 1}

There are $1+3+5+...+2003=1002^{2}=1004004$ small equilateral triangles. 

 Each small equilateral triangle needs $3$ toothpicks to make it. 

 But, each toothpick that isn't one of the $1002\cdot3=3006$ toothpicks on the outside of the large equilateral triangle is a side for $2$ small equilateral triangles. 

 So, the number of toothpicks on the inside of the large equilateral triangle is $\frac{1004004\cdot3-3006}{2}=1504503$.

 Therefore the total number of toothpicks is $1504503+3006=\boxed{\mathrm{(C)}\ 1,507,509}$. ~dolphin7

 



\textbf{Solution 2}

We just need to count upward facing triangles because if we exclude the downward-facing triangles, we won't be overcounting any toothpicks. The first row of triangles has $1$ upward-facing triangle, the second row has $2$ upward-facing triangles, the third row has $3$ upward-facing triangles, and so on having $n$ upward-facing triangles in the $n^\text{th}$ row. The last row with $2003$ small triangles has $1002$ upward-facing triangles. By \href{https://artofproblemsolving.com/wiki/index.php?title=Triangular\_number}{Gauss's formula}, the number of the upward-facing triangles in the entire triangle are now $\frac{1002\times1003}{2}$, meaning that the number of toothpicks are $\frac{1002\times1003}{2}\times3$, or $\boxed{\mathrm{(C)}\ 1,507,509}$.

 \textbf{Note} You don't have to calculate the value of $\frac{1002\times1003}{2}\times3$, and you can use units digits to find the answer easily. The units digit of $1002\times1003$ is $6$, and has a unit digit of $3$ after being divided by $2$. Then this is multiplied by $3$, now the final number ending with a $9$. This leaves only one answer choice possible, which is $\boxed{\mathrm{(C)}\ 1,507,509}$.

 



\textbf{Solution 3}

Test out some smaller cases first.

 When there is just $1$ equilateral triangle in the base, you need $3$ toothpicks.When there are $3$ equilateral triangles in the base, you need $9$ toothpicks in all.When there are $5$ equilateral triangles in the base, you need $18$ toothpicks in all.When there are $7$ equilateral triangles in the base, you need $30$ toothpicks in all.

 Taking the finite differences, we get $6, 9, 12.$ It forms a linear sequence. This means the original numbers $(3, 9, 18, 30)$ form a quadratic.

 Let the quadratic be $y = ax^2 + bx + c$ where $y = 2 \cdot \text{(number of equilateral triangles in base)} - 1.$

 Then, we have the following points: $(1, 3), (2, 9), (3, 18), (4, 30).$

 We can plug these values into $y = ax^2 + bx + c$, giving: 

 
 \[a + b + c = 3, 4a + b + c = 9, 9a + 3b + c = 18.\]

 Solving gives $a = b = 1.5, c = 0.$ So 
 \[y = 1.5x^2 + 1.5x.\]

 For our problem, we need $2003$ equilateral triangles in the base. For the quadratic, the corresponding $x$-value would be $\frac{2003 + 1}{2} = 1002$. So our answer is simply: 
 \[1.5 \cdot 1002^2 + 1.5 \cdot 1002 = \boxed{\mathrm{(C)}\ 1,507,509}.\]




\textbf{Solution 4}

We can count the slanted toothpicks and the horizontal toothpicks separately. Since the bottom row has $2003$ triangles, there will be toothpicks between the triangles and on the edges, which makes $2004$ toothpicks on the bottom row. The row above it will have $2$ fewer triangles, so it will have $2$ fewer toothpicks. This continues all the way to the top. Therefore, the number of slanted toothpicks is 
 \[2004 + 2002 + 2000 + \cdots + 4 + 2 = 1002^2 + 1002.\]

 The horizontal toothpicks are easier. For each row, we will only count the toothpicks on the bottom of the row to avoid overcounting. Since the bottom row has $2003$ toothpicks, there will be $1002$ toothpicks below the row and $1001$ toothpicks above the row (which is the bottom of the next row). This continues all the way to the top, so the number of horizontal toothpicks is 
 \[1002 + 1001 + 1000 + \cdots + 2 + 1 = \frac{(1002)(1003)}{2}.\]

 Adding these together, we get our answer of 
 \[1002^2 + 1002 + \frac{(1002)(1003)}{2} = \boxed{\mathrm{(C)}\ 1,507,509}.\]



	\item (\textit{AMC 10 10A 2003}) Sally has five red cards numbered $1$ through $5$ and four blue cards numbered $3$ through $6$. She stacks the cards so that the colors alternate and so that the number on each red card divides evenly into the number on each neighboring blue card. What is the sum of the numbers on the middle three cards? 


 $\mathrm{(A) \ } 8\qquad \mathrm{(B) \ } 9\qquad \mathrm{(C) \ } 10\qquad \mathrm{(D) \ } 11\qquad \mathrm{(E) \ } 12$

 \textbf{Answer: }E



\textbf{Solution}

Let $R_i$ and $B_j$ designate the red card numbered $i$ and the blue card numbered $j$, respectively. 

 $B_5$ is the only blue card that $R_5$ evenly divides, so $R_5$ must be at one end of the stack and $B_5$ must be the card next to it. 

 $R_1$ is the only other red card that evenly divides $B_5$, so $R_1$ must be the other card next to $B_5$. 

 $B_4$ is the only blue card that $R_4$ evenly divides, so $R_4$ must be at one end of the stack and $B_4$ must be the card next to it. 

 $R_2$ is the only other red card that evenly divides $B_4$, so $R_2$ must be the other card next to $B_4$. 

 $R_2$ doesn't evenly divide $B_3$, so $B_3$ must be next to $R_1$, $B_6$ must be next to $R_2$, and $R_3$ must be in the middle. 

 This yields the following arrangement from top to bottom: $\{R_5,B_5,R_1,B_3,R_3,B_6,R_2,B_4,R_4\}$

 Therefore, the sum of the numbers on the middle three cards is $3+3+6=\boxed{\mathrm{(E)}\ 12}$.

 


	\item (\textit{AMC 10 10A 2003}) Let $n$ be a $5$-digit number, and let $q$ and $r$ be the quotient and the remainder, respectively, when $n$ is divided by $100$. For how many values of $n$ is $q+r$ divisible by $11$? 


 $\mathrm{(A) \ } 8180\qquad \mathrm{(B) \ } 8181\qquad \mathrm{(C) \ } 8182\qquad \mathrm{(D) \ } 9000\qquad \mathrm{(E) \ } 9090$

 \textbf{Answer: }B



\textbf{Solution 1}

$11|(q+r)$ implies that $11|(99q+q+r)$ and therefore $11|(100q+r)$, so $11|n$. Then, $n$ can range from $10010$ to $99990$ for a total of $\boxed{8181\Rightarrow \mathrm{(B)}}$ numbers.

 



\textbf{Solution 2}

When a $5$-digit number is divided by $100$, the first $3$ digits become the quotient, $q$, and the last $2$ digits become the remainder, $r$. 

 Therefore, $q$ can be any integer from $100$ to $999$ inclusive, and $r$ can be any integer from $0$ to $99$ inclusive. 

 For each of the $9\cdot10\cdot10=900$ possible values of $q$, there are at least $\left\lfloor \frac{100}{11} \right\rfloor = 9$ possible values of $r$ such that $q+r \equiv 0\pmod{11}$. 

 Since there is $1$ "extra" possible value of $r$ that is congruent to $0\pmod{11}$, each of the $\left\lfloor \frac{900}{11} \right\rfloor = 81$ values of $q$ that are congruent to $0\pmod{11}$ have $1$ more possible value of $r$ such that $q+r \equiv 0\pmod{11}$. Another way to think about it is the number of possible values of q when r, the remainder, is $0$. In this case, q itself has to be a multiple of $11$. $\left\lfloor \frac{999}{11} = 90 \right\rfloor$. Then we'll need to subtract $9$ from $90$ since we only want multiples of $11$ greater than $100$ $(90-9=81)$

 Therefore, the number of possible values of $n$ such that $q+r \equiv 0\pmod{11}$ is $900\cdot9+81\cdot1=8181 \rightarrow B$.

 



\textbf{Solution 3}

Let $n$ equal $\overline{abcde}$, where $a$ through $e$ are digits. Therefore,

 $q=\overline{abc}=100a+10b+c$
$r=\overline{de}=10d+e$

 We now take $q+r\bmod{11}$:

 $q+r=100a+10b+c+10d+e\equiv a-b+c-d+e\equiv 0\bmod{11}$

 The \href{https://artofproblemsolving.com/wiki/index.php?title=Divisibility\_rules\#Divisibility\_Rule\_for\_Divisibility\_rules\#Divisibility\_Rule\_for\_11}{divisor trick for 11} is as follows:

 "Let $n=\overline{a_1a_2a_3\cdots a_x}$ be an $x$ digit integer. If $a_1-a_2+a_3-\cdots +(-1)^{x-1} a_x$ is divisible by $11$, then $n$ is also divisible by $11$."

 Therefore, the five-digit number $n$ is divisible by 11. The 5-digit multiples of 11 range from $910\cdot 11$ to $9090\cdot 11$. There are $8181\Rightarrow \mathrm{(B)}$ divisors of 11 between those inclusive.

 



\textbf{Solution 4}

Since $q$ is a quotient and $r$ is a remainder when $n$ is divided by $100$. So we have $n=100q+r$. Since we are counting choices where $q+r$ is divisible by $11$, we have $n=99q+q+r=99q+11k$ for some $k$. This means that $n$ is the sum of two multiples of $11$ and would thus itself be a multiple of $11$. Then we can count all the five-digit multiples of $11$ as in Solution 2. (This solution is essentially the same as Solution 3, but it does not necessarily involve mods and so could potentially be faster.)

 



\textbf{Solution 5}

Defining $q$ and $r$ in terms of floor functions and $n$ would yield: $q=\left \lfloor \frac{n}{100} \right \rfloor$ and $r=n-100 \left \lfloor \frac{n}{100} \right \rfloor$. Since $q+r \equiv 0\pmod{11}$, $\left \lfloor \frac{n}{100} \right \rfloor + n-100 \left \lfloor \frac{n}{100} \right \rfloor \equiv 0\pmod{11}$. Simplifying gets us $n-99 \left \lfloor \frac{n}{100} \right \rfloor\equiv 0\pmod{11} \rightarrow n \equiv 0\pmod{11}$ ($99 \left \lfloor \frac{n}{100} \right \rfloor\equiv 0\pmod{11}$ is always true since floor function always yields an integer, and 99 is divisible by 11 w/o any remainder). After we come to this conclusion, it becomes easy to solve the rest of the problem ($\left \lfloor \frac{n}{99999} \right \rfloor - \left \lfloor \frac{n}{10000} \right \rfloor$). ~hw21

 


\end{enumerate}